%قبل از شروع مسئله ترجیحاً چیزی ننویسید. اما اگر خواستید بنویسید، در قالب لیستی (فقرات و شمارش) بنویسید.
%برای اضافه کردن نمره، بعدی از شروع‌مسئله، [x] را اضافه کنید و به جای x نمره مورد نظر را قرار دهید. برای زیرمسئله‌ها هم به همین صورت است.
%\شروع‌مسئله[10]
\شروع‌مسئله 
در مرحله‌ٔ آخر شما با استفاده از پیاده‌سازی‌های مراحل قبلتان یک الگوریتم برخورد چندضلعی را به طور کامل پیاده‌سازی می‌کنید. به این صورت که تعدادی چندضلعی در صفحه  به شما داده می‌شود و شما باید برخوردهای بین این چندضلعی‌ها را خروجی دهید. در مراحل قبل کدهایی زدید که می‌توانند به بهینه‌سازی و پیاده‌سازی این الگوریتم به شما کمک کنند. در 
\href{https://en.wikipedia.org/wiki/Collision_detection}{این لینک}
 می‌توانید با روش‌های دیگری، علاوه بر روش‌های دو مرحله‌ٔ قبل، برای بهینه‌سازی الگوریتم این مسأله آشنا شوید. دقت کنید که مجاز به استفاده از کدهای موجود در اینترنت نیستید. از \lr{Python} و \lr{Pypy} می‌توانید استفاده کنید و امکان استفاده از هیچ کتابخانه‌ای را ندارید. همچنین امکان استفاده از روش‌های \lr{multithreading} را ندارید و برنامه‌تان باید به صورت ترتیبی اجرا شود. به ازای یک ورودی برنامه‌تان باید  همیشه جواب ثابتی تولید کند. کد شما باید کاملاً درست کار کند. به این معنا که بر روی تمام تست‌کیس‌ها پاسخ صحیحی را خروجی دهد. نمره‌دهی در این مرحله به صورت رقابتی است. به این معنا که مجموع زمان هر نفر بر روی تست کیس‌ها محاسبه می‌شود و افراد بر اساس این مجموع زمان مرتب می‌شوند. سپس به نفر اول یک نمره تعلق می‌گیرد و نمره بقیۀ افراد به صورت خطی کاهش می‌یابد. دقت کنید به راه‌حل‌های \lr{brute force} نمره‌ای تعلق نمی‌گیرد.
 \\
 \\
 \متن‌سیاه{ورودی و خروجی:}
 شما باید ورودی را از فایل \lr{input.txt} بخوانید و خروجی را در \lr{output.txt} بنویسید. چندضلعی‌ها به صورت مثلث‌بندی شده به ورودی داده می‌شوند. در خط اول تعداد چندضلعی‌ها داده می‌شود. سپس چندضلعی‌ها به ورودی داده می‌شوند. به این صورت که برای هر چندضلعی، در هر خط یک مثلث به صورت 
 $x_1, y_1, x_2, y_2, x_3, y_3$
  داده می‌شود و پس از اتمام مثلث‌های یک چندضلعی \lr{end} می‌آید. مختصات هر مستطیل در بازه‌ی $(-10000, 10000)$ است. در هر خط از خروجی یک برخورد به صورت $(index_1, index_2)$ چاپ می‌شود که $index$ها شماره‌های چندضلعی‌ها هستند و $index_1 < index_2$. شماره‌ی چندضلعی‌ها به ترتیب ورودی دادن آن‌ها است و از $0$ شروع می‌شود. برخوردها باید به ترتیب چاپ شوند.
 \صفحه‌جدید
  
  %\lstinputlisting{source_filename.py}
  \begin{latin}
  \begin{lstlisting}
  # input
  7
  1 0 1 4 3 4
  3 4 3 0 1 0
  end
  5 4 5 7 9 7
  9 7 9 4 5 4
  end
  1 0 1 5 7 5
  7 5 7 0 1 0
  end
  2 2 2 8 12 8
  12 8 12 2 2 2
  end
  1 0 1 10 5 10
  5 10 5 0 1 0
  end
  1 0 1 4 3 4
  3 4 3 0 1 0
  end
  5 4 5 7 9 7
  9 7 9 4 5 4
  end
  # output
  0 2
  0 3
  0 4
  0 5
  1 2
  1 3
  1 6
  2 3
  2 4
  2 5
  2 6
  3 4
  3 5
  3 6
  4 5
  \end{lstlisting}
  \end{latin}

%\متن‌سیاه{سوال عملی اول:}\\
%\متن‌سیاه{دراپ‌تست}\\

%\متن‌سیاه{ورودی:}\\
%در خط اول ورودی، ابتدا k و سپس m داده می‌شود.\\
%\متن‌سیاه{خروجی:}\\
%در تنها سطر خروجی، تعداد آیفون‌های مورد نیاز می‌آید. اگر در هیچ صورتی چنین کاری ممکن نبود ۱- چاپ می‌شود.\\
%\متن‌سیاه{محدودیت‌ها:}\\
%$1\leq k\leq 10^4$

%$m\leq10^9$
%\\
%\textbf{\متن‌سیاه{ورودی و خروجی نمونه}\\
%% تست اول
%\begin{latin}
% \begin{tabularx}{\textwidth}{|X | X|}
% \hline
% \textbf{stdin} &
% \textbf{stdout}\\
% \hline
% 100 100&
% 1\\
% \hline
% 
% \end{tabularx}
%\end{latin}
%
%% تست دوم
%
%\begin{latin}
% \begin{tabularx}{\textwidth}{|X | X|}
% \hline
% \textbf{stdin} &
% \textbf{stdout}\\
% \hline
% 340 20&
% 3\\
% \hline
% 
% \end{tabularx}
%\end{latin}}
% صورت اصلی مسئله را در این‌جا بنویسید.
%اگر مسئله چند قسمت داشت به این صورت عمل کنید:
%\شروع{زیرمسائل}
%\زیرمسئله[0] % مسئلۀ ۱آ
%\زیرمسئله % مسئلۀ ۱ب
%\زیرمسئله % مسئلۀ ۱پ
%\زیرمسئله % مسئلۀ ۱ت
%\پایان{زیرمسائل}




%%%%%%%%%%%%%%%%%%%%%%%%%%%%%%%%%%%%%%%%%%%%%%%%%%%%%
% برخی دستورات مفید برای تمرین‌ها. %
%%%%%%%%%%%%%%%%%%%%%%%%%%%%%%%%%%%%%%%%%%%%%%%%%%%%%

% فرمول‌های نمونه:
%$f(x) = \Theta(x)$
%$T(x, y) \leq \log(x) + 2^y + \binom{2n}{n}$
% {\tt code\_function}


% لیست موارد بدون شماره‌گذاری:
%\شروع{فقرات}
% \فقره مورد اول 
% \شروع{فقرات}
% \فقره مورد اول داخلی 
% \شروع{فقرات}
% \فقره اولین مورد خیلی داخلی 
% \فقره دومین مورد خیلی داخلی 
% \پایان{فقرات}
% \فقره مورد دوم داخلی 
% \پایان{فقرات}
% \فقره مورد دوم 
%\پایان{فقرات}

% لیست موارد شماره‌گذاری شده:
%\شروع{شمارش}
% \فقره یک مورد 
% \فقره یک مورد دیگر 
% \فقره باز هم یک مورد دیگر
%\پایان{شمارش}

% برای نوشتن یک فرمول چند خطی که یک شماره بخوره. برای حالت‌های دیگر به سایت زیر مراجعه کنید:
% https://www.overleaf.com/learn/latex/Aligning_equations_with_amsmath
%\begin{align} 
% \begin{split}
% (x+y)^3 &= (x+y)^2(x+y) \\
% &= (x^2+2xy+y^2)(x+y) \\
% &= (x^3+2x^2y+xy^2) + (x^2y+2xy^2+y^3) \\
% &= x^3+3x^2y+3xy^2+y^3
% \end{split} 
%\end{align}

% نوشتن ماتریس و آرایه دوبعدی
%\begin{align}
% A = 
% \begin{bmatrix}
% A_{11} & A_{21} \\
% A_{21} & A_{22}
% \end{bmatrix}
%\end{align}

% اضافه کردن تصویر با شرح. برای اضافه کردن تصویر بدون شرح، فقط دستور درج‌تصویر را استفاده کنید.
%\شروع{شکل}[h!]
%\تنظیم‌ازوسط
% \درج‌تصویر[پهنا=0.5\پهنای‌متن]{img.jpg}
%\شرح{شرح تصویر}
%\برچسب{تصویر:۱}
%\پایان{شکل}
%\پایان{زیرمسائل}

%کد زیر را ببینید.
%%\lstinputlisting{source_filename.py}
%\begin{latin}
%\begin{lstlisting}
%# Import the modules
%import sys
%import random
%
%ans = True
%
%while ans:
%question = raw_input("Ask the magic 8 ball a question: (press enter to quit) ")
%
%answers = random.randint(1,8)
%
%if question == "":
%sys.exit()
%
%elif answers == 1:
%print "It is certain"
%
%elif answers == 2:
%print "Outlook good"
%
%elif answers == 3:
%print "You may rely on it"
%
%elif answers == 4:
%print "Ask again later"
%
%elif answers == 5:
%print "Concentrate and ask again"
%
%elif answers == 6:
%print "Reply hazy, try again"
%
%elif answers == 7:
%print "My reply is no"
%
%elif answers == 8:
%print "My sources say no"
%\end{lstlisting}
%\end{latin}